\section{Conclusion}
\label{sec:conclusion}
In this note we have demonstrated exact equivalence between the Schur index for
$\mathcal{N}=4$ SYM with group $SO(2n+1)$ and that of the same theory with gauge
group $Sp(2n)$. We did so via two formulations, which can be respectively
thought of as a ``bottom-up'' and ``top-down'' argument. The first was a proof
by construction applied to the narrower unflavoured case. The second was for the
full flavoured case and worked by taking a known theorem which holds for a
wide class of elliptic hypergeometric integrals and showing that it specialises
to the Schur index. The more interesting line of inquiry is the first argument
applied to the unflavoured case. This provides a strategy that may hold for
other dualities one might want to study. The argument for the flavoured case, on
the other hand, while more general and more concise, rests on a theorem in the
literature which is, subjectively, not as illuminating.

We are confident that one can approach the flavoured index in a similar manner
to our unflavoured approach. However we view the resulting argument as not
particularly insightful while being considerably more complex. We instead chose
to present the much simpler result for the unflavoured case to contrast with the
modestly complex argument for the flavoured case, demonstrating that this
question can be answered directly or via consulting known results in the
literature.

We close with a few lines of further inquiry. Theorem~\ref{thm:rains} is a
consequence of the braid identity in elliptic hypergeometric integral theory~\cite{Rains:2018kernel}.
It would be interesting to investigate whether this directly leads to the
kernel factorisation in the unflavoured case, and thus represents a deeper
relationship between the two sets of indices. Another next step is to study
a refinement of the Schur index such as the Macdonald index and show that $S$-duality holds there as well. Finally, demonstrating
equality of the full superconformal index, resolving the conjecture of
Spiridonov and Vartanov, is another question worth answering. We do not expect
Theorem~\ref{thm:rains} on its own to help here, given it bears little
resemblance to the full superconformal index. However there may still be
lessons to be learned from similar constructions.

\textbf{Acknowledgements.} I would like to thank Carlo Meneghelli for first
drawing my attention to this orbit of problems circa 2018. I would also like to thank
Julian Gould for encouraging me to write this up and share it.
